\section{Anhang}
\subsection{Vertiefung Taylor-Reihe}

Dieser Ansatz erlaubt Taylor-Reihen für rationale Funktionen ohne aufwendige Ableitungen via geom. Reihe.
\begin{itemize}
    \item \textbf{1. Partialbruchzerlegung (PBZ):} Komplexe Bruchfunktion wird in Brüche zerlegt.
    \begin{equation*}
        f(x) = \frac{P(x)}{Q(x)} \xrightarrow{PBZ} \frac{A}{ax+b} + \frac{B}{cx+d}
    \end{equation*}

    \item \textbf{2. Geometrische Reihe:} Terme auf geschlossene Form der geom. Reihe bringen:
    \begin{equation*}
        \frac{1}{1-\cor{q}} = \sum_{n=0}^{\infty} \cor{q}^n
    \end{equation*}

    \item \textbf{3. Nenner-Normierung:} Konstanten ausklammern, um exakt eine $1$ am Anfang zu erhalten, Standardform $\frac{1}{1-q}$.
    \begin{equation*}
        \frac{A}{a+bx} = \frac{A}{a} \cdot \frac{1}{1 - \left(-\frac{b}{a}x\right)} \quad \implies \quad \cor{q} = -\frac{b}{a}x
    \end{equation*}

    \item \textbf{4. Linearität und Vorfaktor:} Konstanten und Vorzeichen in die Summe ziehen. Reihen durch $x^n$ in einer Summe vereinen.

    \item \textbf{5. Konvergenzradius ($\rho$):} Geom. Reihe konvergiert nur für $|q| < 1$. Bei $f$ aus mehreren Reihen jede Bedingung prüfen. Der totale Konvergenzradius ist die Schnittmenge (\textbf{Minimum} der Radien).
\end{itemize}

\subsection{Indexverschiebung in Summen}

Um den Summenindex um eine Konstante $s$ zu verschieben, gilt folgende allgemeine Transformation:

\begin{equation*}
\sum_{n=n_0}^{N} a_n \quad \xrightarrow{k = n + s} \quad \sum_{k=n_0+s}^{N+s} a_{k-s}
\end{equation*}

Die logischen Schritte sind drei:
\begin{enumerate}
    \item \textbf{Neue Variable}: Man definiert $k = n + s$.
    \item \textbf{Neue Grenzen}: unten wird $k_{start} = n_0 + s$ und oben $k_{end} = N + s$.
    \item \textbf{Substitution}: $n$ durch $k$ ausdrücken ($n = k - s$) und im Argument ersetzen.
\end{enumerate}

\begin{equation*}
\sum_{n=0}^{\infty} \frac{x^{n+1}}{(n+1)^2} \quad \xrightarrow{k = n + 1} \quad \sum_{k=1}^{\infty} \frac{x^k}{k^2}
\end{equation*}
Da in diesem Fall $n = k - 1$ gilt, wird der Term $(n+1)$ einfach zu $k$.
\subsection{Kontaktordnung $n$}

Um die Kontaktordnung $n$ bei $x_0$ zwischen $f(x)$ und $g(x)$ zu finden, kann man beide in eine \textbf{Taylor-Reihe} entwickeln und die Koeffizienten $a_k$ vergleichen.
Zwei Funktionen haben einen Kontakt der Ordnung $n$, wenn: die ersten $n+1$ Koeffizienten identisch sind.
\textbf{Achtung:} 
\begin{itemize}
    \item Nullkoeffizienten ($a_k = 0$) sind wichtig und müssen beim Vergleich mitgezählt werden.
    \item Der erste abweichende Koeffizient bestimmt, wo der Kontakt "bricht".
\end{itemize}

\subsection{Singularitätstest Ergänzung}\label{chap:sing_erg}
Oft wird die Gleichung durch Terme mit der abhängigen Variable dividiert. Dies schliesst implizit Lösungen aus, die diese nullsetzen.
\subsubsection{Allgemeiner Ablauf}

\begin{enumerate}
    \item \textbf{Kritische Punkte:} 
    Finde jeden Term $g(y)$ oder $f(z)$, der bei der Rechnung in den Nenner gesetzt wurde.
    \item \textbf{Kandidatensuche:}
    Setze den Term gleich null, um mögliche konstante oder spezielle Lösungen zu finden: $g(y) = 0 \implies y = c $
    
    \item \textbf{DGL (Einsetzen):}
    Setze Funktion $y$ und Ableitung $y'$ in originale DGL ein.
    
    \item \textbf{Auswertung des Resultats:}
    \begin{itemize}
        \item \textbf{Lösung:} Identität immer wahr (z.B. $0 = 0$), \cgn{die Funktion ist eine gültige Lösung.} 
    
        \item \textbf{Test fehlgeschlagen:} $x$-abhängige Gleichung (z.B. $-x^2 = 0$), die Funktion ist \underline{keine} DGL-Lösung.
\end{itemize}
\end{enumerate}

\subsection{Neutralität des Faltungsintegrals}
Berechnet man die partikuläre Lösung $y_p(x)$ via \textbf{Faltungsintegral} mit unterer Grenze $0$, gilt geometrisch und analytisch immer:
\[ y_p(0) = 0 \quad \text{e} \quad y_p'(0) = 0 \]

Sind die Anfangsbedingungen $\mathbf{x_0 = 0}$, ist $y_p$ bei der Berechnung der Konstanten $A$ und $B$ komplett \textbf{neutral}.
\cor{Man kann also $A$ und $B$ finden, indem man die Bedingungen \textbf{nur} auf den homogenen Teil anwendet}